\gr{\section{Related work}
\label{sec:related}

Touch has entered robot learning through each of the three model families we
compare. This section asks four questions of that literature: which sensors are
used, how each family feeds touch to its model, what makes tactile signals hard
to learn from, and on which tasks touch has been shown to help.

\subsection{Tactile sensors in learned manipulation}

Three kinds of sensor dominate. \emph{Vision-based gel sensors} film the inside
of a soft elastomer pad with a small camera, so contact appears as a change in
its shading or in the motion of markers printed on it; GelSight \citep{gelsight}
and DIGIT \citep{digit} are the common designs, and our rig carries four sensors
of this kind. They give high spatial resolution, but at camera rates (25--30\,Hz)
and inside a bulky fingertip. \emph{Piezoresistive arrays} instead measure normal
pressure on a grid of taxels: 3D-ViTac uses flexible $16\times16$ pads costing
about \$20 each \citep{huang2025vitac}, and Touch in the Wild a $12\times32$ array
in a hand-held gripper \citep{zhu2025touchwild}. They are thin and cheap but see
only pressure, at a few millimetres' resolution. \emph{Multimodal fingertips}
combine channels: PolyTouch adds acoustic and peripheral vision sensing to a
camera-based tactile sensor and lasts at least twenty times longer than
commercial ones \citep{zhao2025polytouch}; Sparsh-X learns from the Digit~360
fingertip's tactile image, contact audio, motion and pressure together
\citep{higuera2025sparshx}; and TacThru makes the elastomer transparent so that one
camera sees both the contact and the scene beyond it \citep{li2026tacthru}.
Other policies read force rather than images: Reactive Diffusion Policy works
with marker-based gel sensors and with the arm's joint torques
\citep{xue2025rdp}.

\subsection{How the three model families use touch}

\textbf{Behaviour cloning.} Most work adds touch as one more input to a
diffusion or transformer policy, and the main design question is how to fuse
it with vision. 3D-ViTac places each taxel at its 3D position, computed from the
arm's kinematics, inside the scene's point cloud, so a point-cloud diffusion
policy needs no fusion module at all; supplying the same readings as a flat
tactile image made one in-hand task worse than using no touch
\citep{huang2025vitac}. ViTacFormer, an ACT descendant, fuses vision and touch by
cross-attention in both directions and adds an auxiliary head that predicts the
next tactile reading \citep{heng2026vitacformer}. Touch in the Wild pretrains a
visuo-tactile encoder on data from a portable gripper and conditions a diffusion
policy on it \citep{zhu2025touchwild}; TacThru-UMI feeds simultaneous visual and
tactile signals to a transformer diffusion policy \citep{li2026tacthru}. Two
works change what touch does rather than how it is encoded: Reactive Diffusion
Policy splits the policy into a slow latent diffusion planner (1--2\,Hz) and a
fast tactile feedback loop \citep{xue2025rdp}, and FARM puts grip force into the
action space, predicting a target force alongside the arm pose and regulating
it in closed loop \citep{helmut2026farm}.

\textbf{Vision-language-action models.} Pretrained VLAs have no tactile input,
and large tactile datasets to retrain them do not exist. VLA-Touch therefore
leaves the VLA frozen and wraps it: a tactile-language model turns gel images
into descriptions of hardness or roughness for a language planner, and a
separate controller refines the VLA's action chunks using tactile force
estimates \citep{bi2026vlatouch}. DreamTacVLA instead finetunes a VLA with
tactile images and a tactile world model that predicts future tactile signals,
so that actions are conditioned on imagined contact \citep{ye2025dreamtacvla}.
Our pi0.5 takes the simplest route, feeding fingertip images into image slots
pretrained on wrist cameras (Section~\ref{sec:models}).

\textbf{World models and world action models.} Predicting touch alongside video
is very recent. Visuo-tactile world models predict future visual and tactile
states and use them to plan, keeping objects in place and obeying the laws of
motion better than vision-only models \citep{higuera2026vtwm}. OmniVTA predicts
short-horizon contact evolution with a two-stream world model and corrects the
resulting actions with a 60\,Hz reflex controller \citep{zheng2026omnivta}.
Dream-Tac is a world action model that jointly predicts actions, future images
and future tactile signals \citep{lou2026dreamtac}. DreamZero and FastWAM, which
we use, have no tactile channel of their own: we tile the fingertip images into
their video frame (Section~\ref{sec:models}), so they predict touch only as
pixels of that frame.

\subsection{What makes touch hard to learn from}

Four difficulties recur. \emph{Touch is sparse in time}: it carries information
only during contact, so averaged over a demonstration it can look unimportant
while being decisive at a few moments; 3D-ViTac finds touch irrelevant to
contact-free stages of its tasks \citep{huang2025vitac}. \emph{Representation
matters more than it seems}: the same signal can help or hurt depending on how
it is presented \citep{huang2025vitac}, and frozen pretrained touch encoders work
well for perception but struggled to drive a closed-loop policy in Sparsh
\citep{higuera2025sparsh}. \emph{Sensors differ}: illumination, markers and gel
geometry vary between devices, which is why Sparsh pretrains across three gel
sensors \citep{higuera2025sparsh}, and why an encoder needing four channels from
one fingertip is tied to that device \citep{higuera2025sparshx}. \emph{Timing}: a
gel sensor at 30 frames per second misses vibrations shorter than a frame, which
motivates audio and inertial channels \citep{higuera2025sparshx,zhao2025polytouch}
and fast feedback loops beside a slow planner \citep{xue2025rdp,zheng2026omnivta}.
Optical sensors add their own artefacts, such as light at the gel's edges, which
is the question our crop experiments ask (Section~\ref{sec:crop}).

\subsection{Where touch has been shown to help}

Touch helps most where the deciding state is invisible: under occlusion, inside
the hand, or in the force of a contact. Reported examples include peg and plug
insertion \citep{heng2026vitacformer,higuera2025sparshx}, test-tube insertion and
pipetting \citep{zhu2025touchwild}, in-hand adjustment and fragile objects such as
eggs \citep{huang2025vitac}, peeling and wiping with a held tool
\citep{xue2025rdp,bi2026vlatouch}, picking grapes without crushing them and
keeping a tool seated \citep{helmut2026farm}, and judging hardness or roughness
\citep{bi2026vlatouch}. 3D-ViTac measures the effect directly: touch added more
success with one camera than with three \citep{huang2025vitac}. Folding cloth with
two arms, our task, is less studied. Its contacts are long grasps of a thin,
compliant material rather than precise insertions, and whether touch helps
there, and in which model family, is the question this report starts to answer.}
